\hypertarget{conclusion}{%
\chapter{Conclusion}\label{conclusion}}

\hypertarget{summary-of-contribution}{%
\section{Summary of Contribution}\label{summary-of-contribution}}

This thesis set out to design, build, and evaluate an end-to-end
decision-support pipeline that predicts the win probability of an
Engineering BU proposal and explains that prediction in a form a
business stakeholder can act on, framed as Design Science Research
\citep{peffers2007design, hevner2004design} and structured internally by
CRISP-DM \citep{chapman2000crispdm}. The delivered contribution has four
parts: a synthetic bid dataset, generated from the Engineering BU's own
documented feature hypotheses and deliberately designed, through noise,
correlation, interaction effects, class imbalance, and a
counter-intuitive encoded pattern, to avoid trivial circularity (Chapter
3); a gradient-boosted classification model, evaluated both on a
held-out split and by cross-validation and found to sit close to the
dataset's own recoverable-signal ceiling (Chapter 5); a SHAP explanation
layer providing both a global feature-importance ranking and per-bid
attribution, including a systematic comparison of that ranking against
the BU's a priori hypotheses (RQ1, Chapter 5); and a two-part
interactive artefact, a results dashboard and a live, exactly-auditable
what-if scorer built on a transparent logistic surrogate, published in
two independently accessible forms (Chapter 4).

Against the three research questions: RQ1 found that the model's
top-ranked features agree with the BU's own top hypotheses, but that the
full ranking correlation is weak (Spearman ρ = 0.242), with the largest
divergences individually explicable through interaction effects,
value-range differences, and correlated signal rather than indicating
pipeline failure. RQ2 found a cross-validated AUC of 0.705, close to the
dataset's logistic-regression ceiling of approximately 0.72, indicating
the pipeline recovers close to the full signal the synthetic data
carries, while leaving open whether a more flexible model's advantage
would be larger on bigger, real data. RQ3 found that the artefact meets
four concrete usability requirements derived from the research question
itself, and that building it surfaced two genuine, unplanned design
findings, a hosting-platform data-shape constraint and the case for a
separate, exactly-interpretable scorer model, that are themselves part
of this thesis's contribution to how this kind of artefact should be
built.

Two deliberate deviations from the thesis's original proposal, a single
gradient-boosted model rather than a Random Forest-versus-XGBoost
comparison, and a web dashboard rather than Power BI, were made explicit
from Chapter 1 onward rather than left for a reader to discover, with
the reasoning for each given in full in Chapter 3 and revisited as
limitations in Chapter 6. Both are treated here not as scope abandoned
but as scope deferred: the pipeline was deliberately built so that both
extensions are straightforward additions to existing, working
infrastructure rather than work that must start from nothing.

\hypertarget{recommendations-for-future-work}{%
\section{Recommendations for Future
Work}\label{recommendations-for-future-work}}

\textbf{Move to real bid data.} This is the single highest-priority next
step and the one every other recommendation below depends on for its
real-world relevance. The Engineering BU should begin recording resolved
bid outcomes against the thirteen-feature structure used in this thesis
(Chapter 3), so that the existing pipeline, the data preparation, model
training, SHAP explanation, and artefact-generation code, can be re-run
on real data as soon as a sufficient history accumulates, with no change
to the pipeline itself beyond pointing it at a new input file. A
realistic minimum sample size to begin with should be judged against the
± 0.052 cross-validation variance this thesis found on 500 synthetic
bids (Chapter 5); a real dataset of a similar or larger size would be a
reasonable first target.

\textbf{Run the originally proposed Random Forest-versus-XGBoost
comparison.} Because the synthetic data, feature engineering, and
evaluation harness built for this thesis already exist and are reusable,
adding Random Forest \citep{breiman2001random} and XGBoost
\citep{chen2016xgboost} as additional candidate models and comparing all
three on the same held-out split and cross-validation folds is a
well-defined, comparatively low-effort extension, and should be the
first modelling change made once real data is available, both because it
restores the thesis's original scope and because the comparison itself
may behave differently on larger, real data than the literature review's
general expectation (Section 2.2) predicts.

\textbf{Confirm whether the per-level relationship-history pattern was
recovered.} Chapter 5 identified this as a specific, known gap in the
current evaluation: adding a SHAP dependence plot or partial-dependence
analysis for the relationship-history feature specifically, showing its
effect at each of the four ordinal levels rather than only its aggregate
importance, would directly test whether the pipeline recovered the
deliberately non-monotonic pattern (prior-poor-delivery scoring worse
than new-client) that Chapter 3 built into the data specifically to test
this.

\textbf{Port the pipeline's output to Power BI.} Appendix B sets out the
concrete steps: the structured results file
\texttt{outputs/model\_results.json} was deliberately built as a
portable, self-contained export rather than in a format tied to the web
dashboard, specifically so this porting work does not require re-running
the modelling or SHAP pipeline, only building a new presentation layer
against existing output. This would restore the original proposal's
delivery mechanism and additionally demonstrate the enterprise
deployment features (row-level security, scheduled refresh) the current
web artefact does not.

\textbf{Add threshold selection and multi-class loss reasons to the
artefact.} Chapter 6 identified both as unresolved: letting a user
adjust the scorer's classification threshold interactively, rather than
fixing it at 0.5, would make the precision/recall trade-off identified
in Chapter 5 directly actionable rather than only reported; and
extending the outcome label from binary Won/Lost to the three-class
scheme Zahid et al. \citeyearpar{zahid2021predicting} report adding
decision-relevant value (no bid, client did not pursue, lost to
competition) would require a labelled real dataset that records which of
these actually occurred, making it a natural second step once real data
collection (the first recommendation above) is under way.

\textbf{Conduct a user study with Engineering BU stakeholders.} Once
real data and a real-decision context exist, evaluating RQ3 directly,
observing BU stakeholders using the dashboard and scorer to inform an
actual bid decision, would replace this thesis's structural, author-led
evaluation (Section 5.6) with independent evidence of the artefact's
practical usability, and would likely surface further, currently
unanticipated design requirements in the way the hosting-platform
constraint and the scorer-model decision did during this thesis's own
development process (Section 5.6).

\hypertarget{closing-statement}{%
\section{Closing Statement}\label{closing-statement}}

This thesis set out to show that a working, explainable
bid-competitiveness pipeline could be designed, built, and evaluated for
Virtuos' Engineering BU, and to do so honestly about what a pipeline
built on synthetic data, and adapted in two deliberate, disclosed ways
from its original proposal, can and cannot claim. The pipeline and
artefact delivered recover a known, encoded signal correctly under
realistic noise, explain their own predictions in a form a non-technical
stakeholder can use, and are built in a shape that makes the path to
real-data deployment, the next and most important step, a continuation
of this work rather than a restart of it.
